\begin{proposition}[Cheeger--Stokes 对偶与 \texorpdfstring{$L^\infty$}{L∞} 最优边界探针]\label{prop:spg-cheeger-stokes-duality}
设 \(\Omega\subset\RR^n\) 为有界 Lipschitz 域，取标准体积形式 \(\mathrm{vol}:=\dd x_1\wedge\cdots\wedge \dd x_n\)。
定义 Stokes--Cheeger 常数
\[
\mathcal{C}_\infty(\Omega)
:=\inf\Bigl\{\ \|\eta\|_\infty:\ \eta\in \Omega^{n-1}(\Omega),\ \dd\eta=\mathrm{vol}\ \Bigr\},
\]
其中 \(\|\eta\|_\infty\) 表示系数的 \(L^\infty\) 范数（与本节 \(\Omega^k_{\mathrm{sm}}(I^n)\) 的系数范数一致）。
再定义 Cheeger 常数
\[
\mathfrak h(\Omega)
:=\inf\left\{\ \frac{\mathcal H^{n-1}(\partial^\ast E\cap \Omega)}{|E|}:\ E\subset\Omega,\ 0<|E|<\infty,\ E\ \text{为有限周长集}\right\},
\]
其中 \(|E|\) 为 Lebesgue 体积、\(\partial^\ast E\) 为约化边界。
若存在 Cheeger 极小集 \(E_\ast\) 并且其正则性足以保证存在“校准”向量场 \(V_\ast\in L^\infty(\Omega;\RR^n)\) 满足
\[
\mathrm{div}\,V_\ast=1\ \text{于 }\Omega\ (\text{分布意义}),\qquad
|V_\ast|\le \mathfrak h(\Omega)^{-1}\ \text{a.e.},\qquad
V_\ast\cdot\nu_{E_\ast}=\mathfrak h(\Omega)^{-1}\ \text{a.e. 于 }\partial^\ast E_\ast\cap\Omega,
\]
则有严格等式
\[
\boxed{\ \mathcal{C}_\infty(\Omega)=\mathfrak h(\Omega)^{-1}\ }.
\]
\end{proposition}
\begin{proof}
将 \((n{-}1)\)-形式写为 \(\eta=\iota_V\mathrm{vol}\)，则 \(\dd\eta=(\mathrm{div}\,V)\mathrm{vol}\)，故约束 \(\dd\eta=\mathrm{vol}\) 等价于 \(\mathrm{div}\,V=1\)。
并且在固定的 Euclid 度量下，\(\|\eta\|_\infty\) 与 \(\|V\|_\infty\) 只差一个维数常数；在本文的系数范数约定下可直接以 \(\|\eta\|_\infty=\|V\|_\infty\) 读出。
\par
对任意有限周长集 \(E\subset\Omega\)，有限周长版本的散度定理给出
\[
|E|=\int_E \mathrm{div}\,V\,\dd x=\int_{\partial^\ast E\cap\Omega}V\cdot\nu_E\,\dd\mathcal H^{n-1}
\le \|V\|_\infty\,\mathcal H^{n-1}(\partial^\ast E\cap\Omega).
\]
故 \(\|V\|_\infty\ge |E|/\mathcal H^{n-1}(\partial^\ast E\cap\Omega)\)。
对一切 \(E\) 取上确界得到 \(\|V\|_\infty\ge \mathfrak h(\Omega)^{-1}\)，再对可行 \(V\) 取下确界得到
\(
\mathcal{C}_\infty(\Omega)\ge \mathfrak h(\Omega)^{-1}
\)。
\par
另一方面，若存在上述校准向量场 \(V_\ast\)，则 \(\eta_\ast:=\iota_{V_\ast}\mathrm{vol}\) 满足 \(\dd\eta_\ast=\mathrm{vol}\) 且 \(\|\eta_\ast\|_\infty\le \mathfrak h(\Omega)^{-1}\)，从而
\(
\mathcal{C}_\infty(\Omega)\le \mathfrak h(\Omega)^{-1}
\)。
合并两侧即得结论。证毕。
\end{proof}

\begin{theorem}[超立方体上 \(L^\infty\) 极小原始的边界迹刚性]\label{thm:spg-cube-linfty-primitive-boundary-rigidity}
令 \(k\ge 1\)，取定向立方体 \(I^k=[0,1]^k\) 与常值体积形式
\[
\omega_0:=\dd x_1\wedge\cdots\wedge \dd x_k.
\]
设 \(\eta\in\Omega^{k-1}(I^k)\) 满足 \(\dd\eta=\omega_0\) 且
\[
\|\eta\|_\infty=\frac{1}{2k}.
\]
对每个共维一面 \(F_{j,\varepsilon}:=\{x_j=\varepsilon\}\)（\(\varepsilon\in\{0,1\}\)）记包含映射为 \(\iota_{j,\varepsilon}:F_{j,\varepsilon}\hookrightarrow I^k\)。
则对每个 \(j\) 与 \(\varepsilon\) 都有边界迹刚性：在 \(F_{j,\varepsilon}\) 上几乎处处成立
\[
\boxed{\;
\iota_{j,\varepsilon}^\ast\eta
=(-1)^{j-1}\frac{\varepsilon-\tfrac12}{k}\,
\dd x_1\wedge\cdots\wedge\widehat{\dd x_j}\wedge\cdots\wedge \dd x_k\; }.
\]
从而所有满足 \(\dd\eta=\omega_0\) 且 \(\|\eta\|_\infty=\frac{1}{2k}\) 的极小原始在 \(\partial I^k\) 上具有同一条边界迹。
\end{theorem}
\begin{proof}
由带角点流形的广义 Stokes 定理，
\[
\int_{I^k}\omega_0=\int_{I^k}\dd\eta=\int_{\partial I^k}\eta.
\]
并且 \(\int_{I^k}\omega_0=1\)，故
\begin{equation}\label{eq:spg-cube-linfty-primitive-stokes-identity}
1=\int_{\partial I^k}\eta
=\sum_{j=1}^k\left(\int_{F_{j,1}}\eta-\int_{F_{j,0}}\eta\right).
\end{equation}
对每个 \(F_{j,\varepsilon}\)，由 \(\mathrm{vol}(F_{j,\varepsilon})=1\) 与 \(\|\eta\|_\infty=\frac{1}{2k}\) 得
\[
\left|\int_{F_{j,\varepsilon}}\eta\right|
\le \int_{F_{j,\varepsilon}}|\eta|
\le \frac{1}{2k}.
\]
将其代入 \eqref{eq:spg-cube-linfty-primitive-stokes-identity} 并用三角不等式得到
\[
1
\le
\sum_{j=1}^k\left(\left|\int_{F_{j,1}}\eta\right|+\left|\int_{F_{j,0}}\eta\right|\right)
\le 2k\cdot\frac{1}{2k}=1,
\]
故上述两处不等式必须处处取等。
特别地，对每个 \(F_{j,\varepsilon}\) 有
\(
\left|\int_{F_{j,\varepsilon}}\eta\right|=\frac{1}{2k}
\)
且 \(|\eta|=\frac{1}{2k}\) 在 \(F_{j,\varepsilon}\) 上几乎处处成立。
\par
另一方面，在 \(F_{j,\varepsilon}\) 上 \((k-1)\)-形式空间局部一维：
存在可测函数 \(f_{j,\varepsilon}\) 使
\[
\iota_{j,\varepsilon}^\ast\eta
=f_{j,\varepsilon}\,\dd x_1\wedge\cdots\wedge\widehat{\dd x_j}\wedge\cdots\wedge \dd x_k.
\]
由系数范数界可读出 \(|f_{j,\varepsilon}|\le \frac{1}{2k}\) 几乎处处；
再由 \(\mathrm{vol}(F_{j,\varepsilon})=1\) 与 \(\left|\int_{F_{j,\varepsilon}}\eta\right|=\frac{1}{2k}\) 得 \(f_{j,\varepsilon}=\pm\frac{1}{2k}\) 几乎处处。
最后，\eqref{eq:spg-cube-linfty-primitive-stokes-identity} 的取等迫使
\(
\int_{F_{j,1}}\eta=\frac{1}{2k}
\)
与
\(
\int_{F_{j,0}}\eta=-\frac{1}{2k}
\)
相对于 \(\partial I^k=\sum_j(F_{j,1}-F_{j,0})\) 的诱导定向同时成立。
将诱导定向与坐标楔积形式的关系吸收为因子 \((-1)^{j-1}\)，即得所示边界迹恒等式。证毕。
\end{proof}

\begin{corollary}[一个实现极小常数的显式原始]\label{cor:spg-cube-linfty-primitive-explicit-minimizer}
令
\[
\eta_\star
:=
\frac1k\sum_{j=1}^k(-1)^{j-1}\Bigl(x_j-\tfrac12\Bigr)\,
\dd x_1\wedge\cdots\wedge\widehat{\dd x_j}\wedge\cdots\wedge \dd x_k.
\]
则 \(\dd\eta_\star=\omega_0\)，且 \(\|\eta_\star\|_\infty=\frac{1}{2k}\)；并且 \(\eta_\star\) 的边界迹满足定理 \ref{thm:spg-cube-linfty-primitive-boundary-rigidity} 的刚性公式。
\end{corollary}
\begin{proof}
直接计算
\[
\dd\eta_\star
=\frac1k\sum_{j=1}^k(-1)^{j-1}\dd x_j\wedge
\dd x_1\wedge\cdots\wedge\widehat{\dd x_j}\wedge\cdots\wedge \dd x_k
=\omega_0.
\]
并且 \(|x_j-\tfrac12|\le\tfrac12\) 给出 \(\|\eta_\star\|_\infty=\frac{1}{2k}\)。
边界迹由把 \(x_j\) 限制到 \(0\) 或 \(1\) 直接得到。证毕。
\end{proof}

